\section{Sintassi DL delle Classi Definite}
\label{sec:app_dl}

Di seguito si riporta la sintassi DL compatta di tutte le 17 classi definite
dell'ontologia. La notazione $\sqcap$ indica intersezione ($\mathtt{owl:intersectionOf}$),
$\exists$ indica restrizione esistenziale ($\mathtt{owl:someValuesFrom}$),
$\neg$ indica complemento ($\mathtt{owl:complementOf}$), $\leq$ indica
cardinalit\`a massima ($\mathtt{owl:maxCardinality}$).

\vspace{4pt}
\noindent
\begin{minipage}{\textwidth}
\small
\begin{align*}
\text{MultiplayerGame}      &\equiv \text{VideoGame} \sqcap \exists\text{hasGameMode.MultiplayerMode} \\
\text{SinglePlayerGame}     &\equiv \text{VideoGame} \sqcap \exists\text{hasGameMode.SinglePlayerMode} \\
\text{CoopGame}             &\equiv \text{VideoGame} \sqcap \exists\text{hasGameMode.CoopMode} \\
\text{PCGame}               &\equiv \text{VideoGame} \sqcap \exists\text{availableOn.PCPlatform} \\
\text{ConsoleGame}          &\equiv \text{VideoGame} \sqcap \exists\text{availableOn.ConsolePlatform} \\
\text{MobileGame}           &\equiv \text{VideoGame} \sqcap \exists\text{availableOn.MobilePlatform} \\
\text{NintendoGame}         &\equiv \text{VideoGame} \sqcap \exists\text{availableOn.NintendoPlatform} \\
\text{PlayStationGame}      &\equiv \text{VideoGame} \sqcap \exists\text{availableOn.PlayStationPlatform} \\
\text{XboxGame}             &\equiv \text{VideoGame} \sqcap \exists\text{availableOn.XboxPlatform} \\
\text{SequelGame}           &\equiv \text{VideoGame} \sqcap \exists\text{sequelOf.VideoGame} \\
\text{StandaloneGame}       &\equiv \text{VideoGame} \sqcap \neg\exists\text{sequelOf.VideoGame} \\
\text{HighlyRatedGame}      &\equiv \text{VideoGame} \sqcap \exists\text{metacriticScore.}\geq_{85} \\
\text{RecentGame}           &\equiv \text{VideoGame} \sqcap \exists\text{releaseDate.}\geq_{2020-01-01} \\
\text{AwardWinningGame}     &\equiv \text{VideoGame} \sqcap \exists\text{wonAward.Award} \\
\text{FranchiseGame}        &\equiv \text{VideoGame} \sqcap \exists\text{belongsTo.Franchise} \\[4pt]
\text{GameWithSingleDeveloper} &\sqsubseteq \text{VideoGame} \sqcap \mathopen{\leq}1\,\text{developedBy} \\
\text{ExclusiveRelease}     &\sqsubseteq \text{VideoGame} \sqcap \mathopen{\leq}1\,\text{availableOn}
\end{align*}
\end{minipage}

\vspace{8pt}
\noindent
Le prime 15 classi utilizzano $\equiv$ (condizioni necessarie e sufficienti,
\texttt{owl:equivalentClass}); le ultime due usano $\sqsubseteq$ (sole condizioni
necessarie, \texttt{rdfs:subClassOf}).
